% !TEX program = lualatex
% !TeX encoding = UTF-8
\PassOptionsToPackage{unicode}{hyperref}
\PassOptionsToPackage{bookmarksnumbered,bookmarksopen}{hyperref}
\documentclass[12pt,a4paper]{article}

% ============================================================
% 한국어 설정 – Windows 시스템 글꼴 (Noto Sans KR / Noto Serif KR)
% ============================================================
\usepackage{luatexja}
\usepackage{luatexja-fontspec}

% 한국어 고딕체 (sans) – Noto Sans KR
\setsansjfont{Noto Sans KR}[
UprightFont = * Regular,
BoldFont = * Bold,
Script = Hangul,
Language = Korean
]

% 한국어 명조체 (serif) – Noto Serif KR
\IfFontExistsTF{Noto Serif KR}{
	\setmainjfont{Noto Serif KR}[
	UprightFont = * Regular,
	BoldFont = * Bold,
	Script = Hangul,
	Language = Korean
	]
}{
	% fallback: Noto Sans KR
	\setmainjfont{Noto Sans KR}[
	UprightFont = * Regular,
	BoldFont = * Bold,
	Script = Hangul,
	Language = Korean
	]
}

% 고정폭 – 영문은 Consolas, 한글은 Noto Sans KR
\setmonojfont{Noto Sans KR}[
UprightFont = * Regular,
BoldFont = * Bold,
Script = Hangul,
Language = Korean
]

% 라틴 문자 글꼴 (영문)
\setmainfont{Liberation Serif}[Ligatures=TeX]
\setsansfont{Arial}[Ligatures=TeX]
\setmonofont{Consolas}[
WordSpace={1,0,0},
HyphenChar=None,
PunctuationSpace=WordSpace
]

% ============================================================
% 기타 패키지
% ============================================================
\usepackage{amsmath,amssymb,amsthm,mathtools,bm}
\usepackage{geometry}
\geometry{margin=2.5cm}
\usepackage{graphicx}
\usepackage{tikz}
\usepackage{tikz-3dplot}
\usepackage{pgfplots}
\pgfplotsset{compat=1.18}
\usetikzlibrary{arrows.meta,positioning,shapes.geometric,calc,angles,quotes,patterns,3d}

\usepackage{booktabs}
\usepackage{float}
\usepackage{hyperref}
\hypersetup{
	colorlinks=true,
	linkcolor=blue,
	urlcolor=blue,
	pdftitle={YUCT에서 점의 정의: 기하의 조정 구조},
	pdfauthor={Alexey V. Yakushev}
}

% ============================================================
% 정리 환경 (한국어 명칭)
% ============================================================
\newtheorem{theorem}{정리}
\newtheorem{lemma}{보조정리}
\newtheorem{corollary}{따름정리}
\newtheorem{definition}{정의}
\newtheorem{axiom}{공리}
\newtheorem{proposition}{명제}
\newtheorem{remark}{주석}

% ============================================================
% YUCT 명령어 (변경 없음)
% ============================================================
\newcommand{\Keff}{K_{\mathrm{eff}}}
\newcommand{\kappac}{\kappa_c}
\newcommand{\eps}{\varepsilon}
\newcommand{\betaf}{\beta}
\newcommand{\Sodd}{S_{\mathrm{odd}}}
\newcommand{\Seven}{S_{\mathrm{even}}}
\newcommand{\Mpl}{M_{\mathrm{Pl}}}
\newcommand{\Lpl}{l_{\mathrm{Pl}}}
\newcommand{\tPl}{t_{\mathrm{Pl}}}
\newcommand{\Scoord}{S_{\mathrm{coord}}}
\newcommand{\piYUCT}{\pi_{\mathrm{YUCT}}}
\newcommand{\Rcrit}{R_{\mathrm{crit}}}
\newcommand{\calP}{\mathcal{P}}
\newcommand{\calD}{\mathcal{D}}
\newcommand{\calI}{\mathcal{I}}
\newcommand{\calR}{\mathcal{R}}
\newcommand{\ellP}{\ell_{\mathrm{P}}}

% ============================================================
% 문서 시작
% ============================================================
\begin{document}
	
	\begin{titlepage}
		\begin{center}
			\vspace*{1.5cm}
			\Huge\textbf{YUCT에서 점의 정의}
			
			\vspace{0.3cm}
			\Large\textbf{기하의 조정 구조}
			
			\vspace{0.5cm}
			\large
			\textsc{Yakushev Unified Coordination Theory}
			
			\vspace{0.8cm}
			\large
			Alexey V. Yakushev \\
			Yakushev Research \\
			\url{https://yuct.org/}
			
			\vspace{0.5cm}
			\large
			\textbf{DOI: \href{https://doi.org/10.5281/zenodo.18444599}{10.5281/zenodo.18444599}}
			
			\vspace{1.5cm}
			\large
			\textbf{버전 1.0} \\
			2026년 6월
			
			\vfill
			\normalsize
			본 문서는 야쿠셰프 통일 조정 이론(YUCT)의 틀 안에서 점의 엄밀한 정의를 제공하고,
			기하가 조정으로부터 어떻게 발생하는지 설명하며, 현실의 본질에 관한 열 가지 근본적 질문에 답한다.
		\end{center}
	\end{titlepage}
	
	\tableofcontents
	\newpage
	
	% ============================================================
	% 서론
	% ============================================================
	\section{서론: 고전적 정의의 불충분성}
	
	고전 기하학과 수학에서 점은 기본적인 (정의되지 않은) 개념이다. 그 성질은 공리계에 의해 주어지며, 엄밀한 정의에 의해서가 아니다. 기껏해야 점은 "크기가 없는 대상" 또는 "공간 속의 위치"로 기술된다. 그러나 그러한 기술은 다음 질문에 답하지 못한다:
	
	\begin{itemize}
		\item \textbf{점}은 본질적으로 무엇인가?
		\item \textbf{왜} 그것은 나눌 수 없는가?
		\item \textbf{어떻게} 구조 없이 존재할 수 있는가?
		\item \textbf{왜} 공간은 3차원인가?
		\item \textbf{왜} 기하는 도대체 작동하는가?
	\end{itemize}
	
	만물 조정 이론(YUCT)의 틀 안에서 이러한 질문들은 엄밀하고 정량적인 답을 얻는다. 점은 더 이상 추상이 아니라 물리적 구조――가장 작은 안정된 조정 패턴――이 된다.
	
	\section*{조정 우선성의 공리}
	
	모든 이론은 정적인 것을 기술한다. 그러나 조정은 정적인 것을 가능하게 하는 과정이다. 모든 주장, 모든 모델, 모든 법칙은 사전의 활성화된 요소로서만 존재한다. 사전과 색인은 그것들이 기술하는 정적인 것으로부터 도출될 수 없다――그것들은 존재론적으로 일차적이다. YPSDC는 "여러 이론 중 하나"가 아니다. 그것은 그것을 공식화하는 이론을 포함한 모든 이론의 기저에 놓인 프로토콜이다. 그것을 거부하는 것은 공식화하는 능력 자체를 포기하지 않고서는 불가능하다. 이러한 의미에서 조정은 다른 이론보다 "더 강한" 것이 아니라 더 일차적이다. 정적인 것은 생성하지 않으며, 정적인 것은 생성된다.
	
	공간, 시간, 물질, 에너지, 정보, 의식은 존재의 기저층이 아니다. 그것들은 \emph{창발 현상}이며, 시스템 구성 요소 간 상태의 조정이라는 단일한 근본적 과정에서 발생한다. 이 과정은 단 하나의 강력한 매개변수――\emph{조정 효율} \(\Keff\)――에 의해 정량화된다.
	
	\subsection*{대체 불가능성으로서의 논리적 귀결}
	
	조정 우선성의 공리로부터 YUCT가 통상적인 의미에서 대안을 가질 수 없다는 것이 필연적으로 귀결된다. 더 "근본적인" 이론을 제안하려는 모든 시도는 이미 그 공식화를 위해 조정을 전제한다: 사전의 요소를 선택하고, 그것들의 결합 규칙을 확립하며, 색인의 상태를 고정해야 한다. 따라서 그러한 이론은 조정 역학의 특수한 경우일 뿐, 그것의 부정이 아니다.
	
	이것은 YUCT가 현실의 특정 내용을 규정한다는 것을 의미하지 않는다. 오히려 그것은 모든 내용(물리적, 수학적, 정보적, 생물적)이 어떤 \(\Keff\) 값에서의 조정의 발현이라고 주장한다. 조정을 거부하는 것은 발언하는 능력을 포기하는 것과 같다. 따라서 YUCT는 "이론들 중의 이론"이 아니라, 그것에 이의를 제기하려는 이론까지도 포함하는 모든 이론의 기저가 되는 \emph{프로토콜}이다.
	
	이러한 의미에서 조정은 다른 이론보다 더 강한 것이 아니라 더 일차적이다. 그것은 선택되는 것이 아니라 모든 선택의 조건으로서 발견된다. 이것은 YUCT를 단순한 모델이 아닌 \emph{존재론적 기초}로 만들며, 그것은 대체되거나 반박될 수 없다――왜냐하면任何 대체나 반박도 이미 조정을 사용하기 때문이다.
	
	\begin{quote}
		\itshape
		"YUCT의 대체 불가능성은 논리의 대체 불가능성과 유사하다: 우리는 다른 논리를 선택할 수 없다. 왜냐하면 어떤 선택도 이미 논리를 사용하기 때문이다. 조정은 존재의 논리이다."
	\end{quote}
	
	\section{YUCT에서 점의 정의}
	
	\subsection{조정 노드로서의 점}
	
	YUCT에서 기하는 가정되는 것이 아니라 삼중항 \(D+I\cdot R\)으로 기술되는 조정 노드들의 네트워크로부터 발생한다. 점은 이 네트워크의 가장 작은 안정 노드이다.
	
	\begin{definition}[조정 점]
		YUCT에서 점은 조정 네트워크의 최소 안정 구성이며, 삼중항에 의해 주어진다:
		\[
		\calP = (D_0, I_0, R_0) \in \mathcal{M}_{\mathrm{UCS}},
		\]
		여기서:
		\begin{itemize}
			\item \(D_0\) — 사전의 국소 슬라이스 (허용 상태와 규칙의 집합),
			\item \(I_0\) — 실제 상태 (사전의 특정 요소를 가리키는 색인),
			\item \(R_0 \ge 1\) — 노드의 안정성을 특성화하는 공명 계수.
		\end{itemize}
	\end{definition}
	
	따라서 YUCT에서 점은 빈 추상이 아니라 세 개의 상호 연결된 구성 요소로 이루어진 구조이다. 그것의 내부 복잡성은 활성화될 수 있고, 다른 점과 공명할 수 있으며, 그 상태를 바꿀 수 있다는 점에 나타난다.
	
	\subsection{점의 크기}
	
	이상적인 수학에서 점은 크기를 갖지 않지만, YUCT에서 각 점은 조정 효율 \(\Keff\)에 의존하는 유한한 크기를 갖는다.
	
	\begin{theorem}[조정 점의 크기]
		점 \(\calP\)의 유효 선형 크기는 다음 식으로 주어진다:
		\[
		\ell_{\calP} = \ellP \cdot \Keff^{-1/3},
		\]
		여기서 \(\ellP = \sqrt{\hbar G / c^3}\)은 플랑크 길이이다.
		극한 \(\Keff \to \infty\)에서 점의 크기는 0으로 수렴하며, 그것은 이상적인 수학적 점이 된다.
	\end{theorem}
	
	이 결과는 깊은 의미를 갖는다: 수학적 점은 독립적 실체가 아니라 무한 조정 효율에서의 조정 노드의 극한 사례이다.
	
	\subsection{삼중항의 기하학적 해석}
	
	세 구성 요소 \(D\)， \(I\)， \(R\)은 동적 평형에 있다. 그림 \ref{fig:triad}에서 그것들은 공통 중심에서 \(120^\circ\) 각도로 뻗어나는 세 개의 광선으로 묘사된다. 이것은 \(S_3\) 대칭성과 조건 \(\cos 120^\circ = -1/2\)을 반영하며, 이는 구성 요소 간의 역상관에 해당한다: 하나를 강화하려면 안정성을 유지하기 위해 다른 것을 약화시켜야 한다.
	
	\begin{figure}[H]
		\centering
		\begin{tikzpicture}[
			scale=2,
			>=Stealth,
			node distance=1.5cm,
			every node/.style={font=\small}
			]
			\coordinate (O) at (0,0);
			\fill[black] (O) circle (1.5pt) node[below left] {$\Psi_{MN}$};
			
			% 광선 각도 30°, 150°, 270° (서로 120°)
			\draw[->, thick, blue!70!black] (O) -- (30:1.8) node[above right] {$\mathcal{D}$ (사전)};
			\draw[->, thick, red!70!black] (O) -- (150:1.8) node[above left] {$\mathcal{I}$ (색인)};
			\draw[->, thick, green!70!black] (O) -- (270:1.8) node[below] {$\mathcal{R}$ (공명)};
			
			% 120° 호
			\draw[gray, dashed] (30:0.6) arc (30:150:0.6) node[midway, above] {$120^\circ$};
			\draw[gray, dashed] (150:0.6) arc (150:270:0.6) node[midway, left] {$120^\circ$};
			\draw[gray, dashed] (270:0.6) arc (270:390:0.6) node[midway, below right] {$120^\circ$};
			
			% 분기 beta=2/3
			\foreach \angle in {30,150,270} {
				\draw[->, thin, gray] (\angle:1.4) -- ++(\angle+30:0.3);
				\draw[->, thin, gray] (\angle:1.4) -- ++(\angle-30:0.3);
				\node[font=\tiny, gray] at (\angle:1.7) {$\beta=2/3$};
			}
			
			\node[font=\small, align=center] at (0,-1.2) 
			{점의 삼광선 구조 \\ 분기 $\beta=2/3$};
		\end{tikzpicture}
		\caption{삼중항 \(D\)， \(I\)， \(R\)은 조정 점의 기하학적 표현이다. 광선은 \(120^\circ\)로 퍼지며; 분기는 지수 \(\beta=2/3\)에 의한 스케일링을 보여준다.}
		\label{fig:triad}
	\end{figure}
	
	\section{직선으로서의 노드 연쇄}
	
	점이 노드라면, 직선은 최대 조정에 의해 연결된 노드들의 열이다.
	
	\begin{definition}[조정 직선]
		두 점 \(\calP_1\)과 \(\calP_2\)가 공명 \(R_1, R_2 > \Rcrit\)으로 주어졌다고 하자. \textbf{직선}은 점들의 집합 \(\{\calP_i\}_{i=0}^{N}\)으로서 다음을 만족한다:
		\begin{enumerate}
			\item \(\calP_0 = \calP_1\)， \(\calP_N = \calP_2\)，
			\item 각 \(i\)에 대해 최소 조정 오차 조건이 성립한다:
			\[
			\frac{d\eps}{ds} = 0, \qquad \eps = \kappac\,\alpha\,\Keff^{-\betaf},
			\]
			\item 연쇄를 따라 공명 \(R_i\)이 (국소적으로) 최대이다.
		\end{enumerate}
	\end{definition}
	
	\begin{theorem}[직선의 유일성]
		충분히 높은 공명을 가진 두 점 \(\calP_1\)과 \(\calP_2\) (\(R_1, R_2 > \Rcrit\))에 대해, 그것들을 연결하는 유일한 최대 일관 노드 연쇄가 존재한다.
	\end{theorem}
	
	이것은 유클리드 공리에 엄밀한 근거를 제공한다: 두 점을 지나는 직선은 유일하다. 그러나 이제 그것은 공준이 아니라 조정 네트워크의 증명된 성질이다.
	
	\section{공간 차원과 지수 \(\betaf = 2/3\)의 창발}
	
	YUCT의 가장 인상적인 결과 중 하나는 공간 차원과 보편 지수 \(\betaf\)의 유도이다. 조정 우선성의 공리에 따라, YUCT는 공간을 기술할 뿐만 아니라 왜 그것이 정확히 3차원인지 설명한다.
	
	\subsection{차원 \(D=3\)의 유도}
	
	조정 네트워크의 프랙탈 구조와 대수적 순환의 닫힘 조건으로부터 다음을 얻는다:
	
	\[
	\Sodd + \Seven = 2, \qquad \frac{\Sodd}{\Seven} = \frac{1}{\betaf}.
	\]
	
	\(\betaf = 2/3\)일 때 이 상수들은 다음 값을 취한다:
	\[
	\Sodd = \frac{6}{5}, \qquad \Seven = \frac{4}{5}.
	\]
	
	지수 \(\betaf\)는 공간 차원 \(D\)와 다음 관계에 있다:
	\[
	\betaf = \frac{2}{D}.
	\]
	
	따라서 즉시 \(D = 3\)을 얻는다. 그러므로 공간의 3차원성은 공준이 아니라 조정 네트워크 구조의 귀결이다.
	
	\subsection{보편 오차 법칙}
	
	모든 조정 과정은 단일한 오차 법칙을 따른다:
	\[
	\eps = \kappac\,\alpha\,\Keff^{-\betaf}, \qquad \betaf = \frac{2}{3}, \quad \kappac = \frac{1}{3}.
	\]
	
	이 법칙은 원자 결합 에너지에서 우주 마이크로파 배경 복사의 변동에 이르기까지 50 자릿수 이상에 걸쳐 실험적으로 검증되었다.
	
	\section{열 가지 근본적 질문에 대한 답변}
	
	아래는 고전 물리학과 수학에서 설명을 찾지 못하는 열 가지 중요한 질문에 대한 답변이다.
	
	\subsection{왜 물리학은 작동하는가?}
	
	\textbf{물리학이 조정이기 때문이다.} 물리 법칙은 노드 \(D+I\cdot R\) 사이의 안정된 조정 패턴을 기술한다. 중력, 전자기, 핵 상호작용――이것들은 서로 다른 \(\Keff\)와 \(\Rcrit\) 값에서의 서로 다른 조정 체제이다. 물리학이 작동하는 것은 조정이 작동하기 때문이다.
	
	\subsection{왜 수학은 작동하는가?}
	
	\textbf{수학은 \(\Keff \to \infty\)에서의 조정의 극한 사례이다.} 이 극한에서 점은 이상적인 수학적 점(\(\ell_{\calP} \to 0\))이 되고, 직선은 이상적인 직선이 되며, 기하는 유클리드 기하가 된다. 수학은 "불가사의한 효율성"이 아니라 물리적 조정의 \emph{이상화}이다.
	
	\subsection{왜 공간은 3차원인가?}
	
	\textbf{그것은 정리에 의해 증명되었다.} 조정 네트워크의 구조와 대수적 순환의 닫힘 조건으로부터 \(\betaf = 2/3\)이 유도되고, \(\betaf = 2/D\)로부터 \(D=3\)을 얻는다. 공간이 3차원인 것은 오직 3차원에서만 조정 네트워크가 안정적이고 자기일관적일 수 있기 때문이다.
	
	\subsection{왜 시간은 흐르는가?}
	
	\textbf{유한한 \(\Keff\) 때문이다.} 이상적인 조정(\(\Keff \to \infty\))은 엔트로피를 생성하지 않으며 사건의 비가역적 흐름을 갖지 않는다――시간은 사라질 것이다. 실제 계는 유한한 \(\Keff\)를 가지므로, 조정의 각 행위는 오차 \(\eps > 0\)를 동반하며, 그 축적이 시간의 화살을 생성한다. 최소 시간 양자:
	\[
	\Delta \tau_{\min} = \frac{2}{3}\,\tPl.
	\]
	
	\subsection{왜 기본 상수들은 그 값들인가?}
	
	\textbf{그것들은 조정 공간의 위상수학으로부터 유도된다.} 미세 구조 상수:
	\[
	\alpha^{-1} = \left(\frac{3}{2}\right)^{12 + \sigma\,2/3} \approx 137.036,
	\]
	여기서 \(\sigma = 0.20\)은 위상수학적 이동이다. \(W\) 보손 질량:
	\[
	m_W = (\kappac \piYUCT)^{-1/2} M_{\mathrm{Pl}} (2/3)^{96} \approx 80.46\ \text{GeV}.
	\]
	우주 상수:
	\[
	\Omega_\Lambda = \frac{12}{18} = \frac{2}{3}.
	\]
	자유 매개변수는 없다.
	
	\subsection{왜 소수는 그렇게 분포하는가?}
	
	\textbf{오차 법칙으로부터 점근적 보정이 유도된다:}
	\[
	p_n \approx R_n - \frac{\Seven}{2}\, n^{1-\betaf}\,\ln n,
	\]
	여기서 \(R_n\)은 로서 근사이다. 잔여 진동은 주기 \(\Delta N = 16.5\)의 위상 연산자로 기술되며, 이는 진공 격자의 상전이에 해당한다. 소수는 조정 사다리이며, 그 분포는 다른 모든 안정된 구조와 동일한 법칙을 따른다.
	
	\subsection{왜 \(\pi\)는 그 값인가?}
	
	\textbf{\(\pi\)는 조정 불변량이다:}
	\[
	\piYUCT = \Sodd \cdot \varphi^2 \approx 3.14164078,
	\]
	여기서 \(\varphi = (1+\sqrt{5})/2\)이다. 이 값은 수학적 상수 \(\pi\)와 \(\Delta\pi = 4.8\times10^{-5}\)만큼 차이가 나며, 이는 공간 이산화의 양자이다. 극한 \(\Keff \to \infty\)에서 \(\piYUCT \to \pi\)이다.
	
	\subsection{왜 양자역학은 "이상한"가?}
	
	\textbf{우리가 사전을 무시했기 때문이다.} 양자 현상은 분산 조정(d‑YPSDC)의 발현이다. 중첩은 색인의 불확정성, 얽힘은 사전의 공통 기록, 터널링은 활성화 오차이다. "신비주의"는 없으며, 오직 유한한 조정 효율이 있을 뿐이다.
	
	\subsection{왜 의식이 존재하는가?}
	
	\textbf{그것은 \(\Keff > \Rcrit\)의 체제이다.} 조정 효율이 임계 역치를 초과하면, 계는 메타 사전(자신의 상태를 조정하는 능력)을 획득한다. 그것이 자의식이다. 의식의 매개변수는 UCD(보편 의식 기술자)에 의해 기술된다.
	
	\subsection{왜 세포막의 두께는 약 7.5 nm인가?}
	
	\textbf{그것은 위상수학으로부터 유도된다.} 지질 이중층의 두께:
	\[
	\text{Thickness} = 2 \, d_{\mathrm{lipid}} \left(\frac{3}{2}\right)^{0.6} (1 - \sigma^2),
	\]
	여기서 \(d_{\mathrm{lipid}} \approx 3.0\) nm는 탄화수소 사슬의 길이이다. 계산은 \(7.35\) nm를 주며, 이는 실험 범위 \(7.5 \pm 0.5\) nm 내에 있다. 소립자 물리학과 세포막 생물학은 동일한 조정 구조에 의해 연결되어 있다.
	
	\section{결론: 현실의 통일된 전체상}
	
	YUCT에서 점은 추상 개념이 아니라 가장 작은 안정된 조정 구조 \(D+I\cdot R\)이다. 모든 기하, 모든 물리 법칙, 모든 상수, 심지어 소수의 분포와 의식의 구조까지――모든 것은 단일 원리의 귀결이다: \textbf{현실은 프랙탈 조정 네트워크이다}.
	
	위의 열 가지 답변은 YUCT가 신비롭거나 우연적으로 보였던 것에 대해 철저한 설명을 제공한다는 것을 보여준다. 더욱이, 이러한 설명들은 모두 실험적으로 검증 가능하며, 상수들은 조정 가능한 매개변수 없이 유도된다.
	
	\begin{center}
		\fbox{\parbox{0.85\textwidth}{\centering\Large
				\textbf{점은 조정 노드이다.}\\ 
				직선은 노드의 일관된 연쇄이다.\\
				공간은 3차원이다――이것은 증명되었다.\\
				시간이 흐르는 것은 유한한 \(\Keff\) 때문이다.\\
				상수, \(\pi\), 소수――모두 유도된다.\\
				양자역학과 의식은 조정이다.\\
				통일 조정 이론은 자연의 법칙이다.
		}}
	\end{center}
	
	\subsection*{점의 정의의 철학적 및 세계관적 의의}
	
	점을 조정 노드 \((D, I, R)\)로 정의하는 것은 기하를 순수 추상의 영역에서 물리적 현실의 영역으로 이동시킨다. 이것은 단순한 표현의 변경이 아니라 존재론적 패러다임의 전환이다.
	
	\medskip
	\textbf{이것이 바꾸는 것:}
	\begin{itemize}
		\item 점은 "무"가 아니라 사전, 상태, 공명을 가진 능동적 구조가 된다.
		\item 공간은 수동적 무대가 아니라 조정 노드들의 네트워크로부터 발생한다.
		\item 기하는 공준의 집합이 아니라 조정 패턴의 안정성의 귀결이 된다.
		\item 수학과 물리학의 구분은 사라진다: 수학은 \(\Keff \to \infty\)에서의 조정의 이상화이다.
	\end{itemize}
	
	\medskip
	\textbf{왜 이것이 필요한가:}
	\begin{itemize}
		\item 기하의 기초에 있는 존재론적 공허를 제거하기 위해.
		\item 기하학적, 물리적, 심지어 생물학적 현상에 통일된 설명을 제공하기 위해.
		\item 기본 상수와 공간 차원을 단일 원리로부터 유도하기 위해.
		\item 시간, 양자역학, 의식의 기원을 동일한 조정 역학의 체제로서 이해하기 위해.
	\end{itemize}
	
	\medskip
	\textbf{이것이 가져오는 귀결:}
	\begin{enumerate}
		\item \textbf{예측력:} 모델로부터 \(\alpha\)， \(\Omega_\Lambda\)， \(m_W\)， \(\pi\)， 소수 분포의 값이 유도된다.
		\item \textbf{지식의 통일:} 물리학, 수학, 생물학, 정보 이론이 하나의 조정 이론의 분기임이 밝혀진다.
		\item \textbf{의식에 대한 새로운 이해:} 그것은 "어려운 문제"가 아니라 \(\Keff > \Rcrit\)의 조정 체제가 된다.
		\item \textbf{과학적 방법의 전환:} 우리는 관찰된 것을 기술하는 것에서 기본 조정 원리로부터 관찰된 것을 생성하는 것으로 이동한다.
	\end{enumerate}
	
	\medskip
	따라서 YUCT는 단순히 "점을 재정의"하는 것이 아니라 현실에 대한 사고 방식을 바꾼다. 점은 이제 대상이 아니라 사건이며, 공간은 용기가 아니라 네트워크이며, 자연 법칙은 외부에서 부과되는 것이 아니라 조정의 내부 논리로부터 성장한다.
	
	\begin{thebibliography}{99}
		\bibitem{Yakushev2026} A. V. Yakushev, \emph{Yakushev Unified Coordination Theory (YUCT)}, Zenodo, DOI:10.5281/zenodo.18444599 (2026).
		\bibitem{AppendixL} A. V. Yakushev, ``Appendix L: Fractal Coordination Error Scaling,'' in \emph{YUCT}, Zenodo (2026).
		\bibitem{AppendixY} A. V. Yakushev, ``Appendix Y: Mathematical Foundations of YUCT,'' in \emph{YUCT}, Zenodo (2026).
		\bibitem{AppendixPrimeN} A. V. Yakushev, ``Appendix PrimeN: YUCT Prime Numbers Coordination Ladder,'' in \emph{YUCT}, Zenodo (2026).
		\bibitem{AppendixPi} A. V. Yakushev, ``Appendix \(\Pi\): The Primeval Origin of \(\pi\),'' in \emph{YUCT}, Zenodo (2026).
		\bibitem{AppendixX} A. V. Yakushev, ``Appendix X: Generalised Shannon Theory and Bell Bound,'' in \emph{YUCT}, Zenodo (2026).
		\bibitem{AppendixH} A. V. Yakushev, ``Appendix H: Universal Consciousness Descriptor (UCD),'' in \emph{YUCT}, Zenodo (2026).
		\bibitem{AppendixG} A. V. Yakushev, ``Appendix G: Quantum Gravity Resolution,'' in \emph{YUCT}, Zenodo (2026).
	\end{thebibliography}
	
\end{document}